\section{Постановка задачи}
Пусть $\mathcal D_s=\{(x_i^s,y_i^s)\}_{i=1}^{N_s}$ --- размеченная исходная выборка, а $\mathcal D_t=\{x_j^t\}_{j=1}^{N_t}$ --- целевая выборка без доступных обучающих меток. Пространство из $C$ классов считается общим. Сеть с параметрами $\theta$ выдаёт в пикселе $u$ вектор вероятностей $p_\theta(x)_u$; итоговая метка выбирается по максимальной вероятности. Распределения изображений двух доменов различаются.

Схема разделения обучающих и проверочных данных приведена на рис.~\ref{fig:uda-protocol}.

\begin{figure}[htbp]
\centering
% Авторская схема. Все подписи набираются шрифтом статьи.
\begingroup
\setlength{\unitlength}{1mm}
\fontsize{8.5}{10}\selectfont
\linethickness{0.4pt}
\begin{picture}(86,61)
  \put(0,49){\framebox(40,12){\shortstack{GTA5: $(x^s,y^s)$\\изображения и маски}}}
  \put(46,49){\framebox(40,12){\shortstack{Cityscapes train: $x^t$\\без обучающих масок}}}
  \put(20,49){\line(0,-1){6}}
  \put(66,49){\line(0,-1){6}}
  \put(20,43){\line(1,0){46}}
  \put(43,43){\vector(0,-1){6}}
  \put(7,25){\framebox(72,12){\shortstack{Обучение общей модели\\$\min_\theta\;L_s+\sum_k\lambda_k L_k^{\mathrm{adapt}}$}}}
  \put(43,25){\vector(0,-1){11}}
  \put(45,19){\makebox(0,0)[l]{параметры $\theta$}}
  \put(0,0){\framebox(25,14){\shortstack{Cityscapes val\\изображения}}}
  \put(25,7){\vector(1,0){6}}
  \put(31,0){\framebox(24,14){\shortstack{Модель\\предсказания}}}
  \put(55,7){\vector(1,0){6}}
  \put(61,0){\framebox(25,14){\shortstack{Сравнение\\с масками val\\mIoU, IoU}}}
\end{picture}
\endgroup

\caption{Постановка GTA5~$\rightarrow$~Cityscapes: обучение без целевых меток и оценка на отдельной проверочной выборке.}
\label{fig:uda-protocol}
\end{figure}

На исходных данных минимизируется перекрёстная энтропия:
\begin{equation}
 L_s(\theta)=-\frac1{N_s}\sum_{i=1}^{N_s}\frac1{|\Omega_i|}
 \sum_{u\in\Omega_i}\log p_{\theta,y^s_{iu}}(x_i^s)_u,
 \label{eq:source}
\end{equation}
где $\Omega_i$ --- множество учитываемых пикселей. Для целевых обучающих изображений такая потеря недоступна. Общая схема вводит дополнительный критерий:
\begin{equation}
 \min_\theta\; L_s(\theta)+\sum_k\lambda_k L_k^{\mathrm{adapt}}(\theta),
 \qquad\lambda_k\geq0.
 \label{eq:total}
\end{equation}
Формула задаёт способ объединения критериев, но не требует одновременного использования всех направлений. Для состязательной схемы добавляется максимизация по параметрам дискриминатора, для транспорта --- внутренняя задача определения соответствий.

Основной сценарий --- GTA5~\cite{gta} как источник синтетических изображений и Cityscapes~\cite{cityscapes} как целевой домен. Целевые обучающие маски не используются при оптимизации. Маски отдельной проверочной выборки нужны для оценки. Если они применяются для выбора гиперпараметров или момента остановки, это должно явно отражаться в протоколе: такой режим отличается от выбора модели без целевой разметки. SYNTHIA может рассматриваться позднее как дополнительный источник; результаты для 13, 16 и 19 классов нельзя смешивать в одной оценке.

Качество оценивается средней по классам мерой пересечения с объединением:
\begin{equation}
 \mathrm{mIoU}=\frac{100}{C}\sum_{c=1}^{C}
 \frac{\mathrm{TP}_c}{\mathrm{TP}_c+\mathrm{FP}_c+\mathrm{FN}_c}.
 \label{eq:miou}
\end{equation}
Числа истинно положительных, ложноположительных и ложноотрицательных пикселей накапливаются по всей проверочной выборке. Дополнительно рассматриваются IoU каждого класса, время обучения и пиковая память ускорителя. Средняя метрика может скрывать ухудшение редких категорий, поэтому одного mIoU недостаточно.
